\documentclass[12pt,aspectratio=169]{beamer}
\input{header_beamer}
\usetheme{Boadilla}
\usecolortheme{beaver}
\setbeamertemplate{page number in head/foot}{}
\setbeamertemplate{navigation symbols}{}
\title{Lecture-09: Moments and $L^p$ spaces}

\begin{document}
\pagestyle{empty}
\maketitle

\begin{frame}{Moments}
Let $X:\Omega\to\R$ be a random variable defined on the probability space $(\Omega,\sF,P)$ with the distribution function $F_X:\R \to [0,1]$. 

\begin{Example}[Absolute value function]<2-> 
For the function $\abs{\cdot}: \R \to \R_+$, 
we can write the inverse image of half open sets $(-\infty, x]$ for any $x \in \R$ as $A_g(x) = g^{-1}(-\infty, x]$. 
It follows that $A_g(x) = \emptyset \in \sB(\R)$ for $x < 0$ and $A_g(x) = [-x, x] \in \sB(\R)$ for $x \in \R_+$. 
Since $g^{-1}(-\infty, x] \in \sB(R)$, it follows that $\abs{\cdot}: \R \to \R_+$ is a Borel measurable function. 
\end{Example}
\end{frame}

\begin{frame}{Moments}
\begin{block}{Lemma}<1->
If $\E\abs{X}$ is finite, then $\E X$ exists and is finite. 
\end{block}
\begin{Proof}<2->
\begin{enumerate}
\item <2-> The function $\abs{\cdot}: \R \to \R$ is a Borel measurable function and hence $\abs{X}$ is a random variable. 
\item <3-> Further $\abs{X} \ge 0$, and hence the expectation $\E\abs{X}$ always exists. 
\item <4-> If $\E\abs{X}$ is finite, it means $\E X_+$ and $\E X_-$ are both finite, 
and hence $\E X = \E X_+ - \E X_-$ is finite as well. 
\end{enumerate}
\end{Proof}
\end{frame}
\begin{frame}{Moments}
\begin{cor}<1->
Let $g: \R \to \R$ be a Borel measurable function. 
If $\E\abs{g(X)}$ is finite, then $\E g(X)$ exists and is finite. 
\end{cor} 

\begin{block}{Exercise : Polynomial function}<2->
For any $k \in \N$, we define functions $g_k: \R \to \R$ such that $g_k: x \mapsto x^k$. 
Show that $g_k$ is Borel measurable for all $k \in \N$. 
\end{block} 
\end{frame}
\begin{frame}{Moments}
\begin{Definition}[Moments]<1->
We define the $k$th \textbf{moment} of the random variable $X$ as $m_k \triangleq \E g_k(X) = \E X^k$. 
First moment $\E X$ is called the \textbf{mean} of the random variable. 
\end{Definition}
\begin{block}{Reminder}<2->
\begin{itemize}
\item <2-> If $\E\abs{X}^k$ is finite, then $m_k$ exists and is finite. 
\item <3-> If $P\set{\abs{X} \le 1} = 1$, then $P\set{\abs{X}^k \le 1} = 1$.  By the monotonicity of expectations $\E\abs{X}^k \le 1$, 
and the moments $m_k$ exist and are finite for all $k \in \N$. 
\end{itemize}
\end{block}
\end{frame}

\begin{frame}{$L^p$ spaces}
\begin{Definition}<1->
A vector space over field $\F$ is denoted by a set $V$ has 
\begin{enumerate}[(a)] 
\item <2-> vector addition $+: V\times V\to V$ that satisfies associativity and commutativity, and has an identity and inverse element for each vector $v\in V$ 
\item <3-> scalar multiplication $\cdot: \F \times V \to V$ that satisfies compatibility of field multiplication, distributivity with vector and field addition, and has an identity element. 
\end{enumerate}
\end{Definition}
\end{frame}

\begin{frame}{$L^p$ spaces}
\begin{block}{Reminder} 
\begin{itemize}
\item <1-> Set of random variables is a vector space over reals, denoted by $V$. %Can verify the associativity and commutativity of addition of random variables. 
\item <2-> The constant function $0$ is the identity random variable, and $-X$ is the additive inverse for vector additions.  
\item <3-> %Can verify the compatibility of scalar multiplication with random variables, 
Can verify the existence of unit scalar $1$, and distributivity of scalar multiplications with respect to field addition and vector addition.  
\end{itemize}
\end{block}
\end{frame}

\begin{frame}{Norm}
\begin{Definition}<1->
For a vector space $V$ over field $\F$, a \textbf{norm} on the vector space is a map $f: V \to \R_+$ 
that satisfies  
\begin{enumerate}
	\item <2-> [\textbf{homogeneity:}] $f(av) = \abs{a}f(v)$ for all $a \in \F$ and $v \in V$,
	\item <3-> [\textbf{sub-additivity:}] $f(v+w) \le f(v) + f(w)$ for all $v,w \in V$, and 
	\item <4-> [\textbf{point-separating:}] $f(v) \ge 0$ for all $v \in V$. 
\end{enumerate}
\end{Definition}
\end{frame}

\begin{frame}{Moments and $L^p$ spaces}
\begin{Definition}<1->
For a probability space $(\Omega,\sF,P)$, and $p \ge 1$, 
we define the set of random variables with finite absolute $p$th moment as the vector space  
$
L^p \triangleq \set{X \in V:(\E\abs{X}^p)^{\frac{1}{p}} < \infty}. 
$
\end{Definition}
\begin{Definition}<2->
We define a function $\norm{}_p: L^p \to \R_+$ defined by $\norm{}_p(X) = \norm{X}_p \triangleq (\E\abs{X}^p)^{\frac{1}{p}}$ for any $X \in L^p$ and real $p \ge 1$. 
\end{Definition}

\begin{block}{Reminder}<3->
For $p=1$, the map $\norm{}_p$ is norm. 
Therefore, $L^1$ is a normed vector space consisting of random variables with bounded absolute mean.
\end{block}
\end{frame}

\begin{frame}{Moments and $L^p$ spaces}
\begin{block}{Reminder}
\begin{itemize}
\item <1-> $\norm{X}_\infty = \sup\set{\abs{X(\omega)}:\omega\in \Omega}$, 
and hence $L^\infty$ is a normed vector space of bounded random variables. 
\item <2-> $\norm{}_p$ is a norm for all $p \in (1, \infty)$, 
and hence $L^p$ is a normed vector space of random variables with bounded $\norm{}_p$ norm. 
In particular, the $L^2$ space consists of random variables with bounded second moment.

\item <3-> If $\E\abs{X}^N$ is finite for some $N \in \N$, 
then $\E\abs{X}^k$ is finite for all $k \in [N]$. 

\item <4-> \EQ{
\abs{X}^k = \abs{X}^k\SetIn{\abs{X} \le 1} + \abs{X}^k\SetIn{\abs{X} > 1} \le 1 + \abs{X}^N. 
}
\item <5-> This implies that $L^N \subseteq L^k$ for all $k \in [N]$. $L^{q} \subseteq L^p$ for any real numbers $1 \le p \le q$. 
\end{itemize}
\end{block}
\end{frame}

\begin{frame}{Central Moments}
\begin{Example}[Shifted polynomial functions]<1->
For any $k \in \N$, we define functions $h_k: \R \to \R$ such that $h_k: x \mapsto (x-m_1)^k$. 
Then, $h_k = g_k(x-m_1) = g_k\circ f$ where $f:\R \to\R$ is defined as $f(x) = x-m_1$ for all $x \in \R$. 
Since $g_k$ and $f$ are measurable, so is $h_k$. 
\end{Example}

\begin{Definition}[Central moments]<2->
\begin{itemize}
\item <2-> Let $X: \Omega \to \R$ be a random variable defined on the probability space $(\Omega, \sF, P)$ with finite first moment $m_1$. 
\item <3-> We define the $k$th \textbf{central moment} of the random variable $X$ as $\sigma_k \triangleq \E h_k(X) = \E (X-m_1)^k$. 
\item <4-> The second central moment $\sigma_2 = \E(X-m_1)^2$ is called the \textbf{variance} of the random variable and denoted by $\sigma^2$. 
\end{itemize}
\end{Definition}
\end{frame}

\begin{frame}{Central moments}
\begin{block}{Lemma}
The first central moment $\sigma_1= \E(X-m_1) = 0$ and the variance $\sigma^2 = \E(X-m_1)^2$ for a random variable $X$ is always non-negative, 
with equality when $X$ is a constant. 
That is, $m_2 \ge m_1^2$ with equality when $X$ is a constant. 
\end{block}
\end{frame}

\begin{frame}{Central moments}
\begin{Proof}
\begin{itemize}
\item <1-> $h_1(X) = X- m_1$ and $h_2(X) = (X-m_1)^2$ are random variables. 
\item <2-> $\sigma_1 = \E h_1(X) = \E X - m_1 = 0$. 
\item <3-> $0 \le \E(X-m_1)^2$.  
\item <4-> $
\sigma^2 = \E X^2 - 2m_1\E X + m_1^2 = m_2 - m_1^2 \ge 0.
$
\end{itemize}
\end{Proof}
\begin{block}{Reminder}<5->
If second moment is finite, then the first moment is finite. 
That is, $L^2 \subseteq L^1$. 
\end{block}
\end{frame}

\begin{frame}{Inequalities}
\begin{block}{Theorem : Markov's inequality}<1->
Let $X: \Omega \to \R$ be a random variable defined on the probability space $(\Omega, \sF, P)$. 
Then, for any monotonically non-decreasing function $f: \R \to \R_+$, we have 
\EQ{
P\set{X \ge \epsilon} \le \frac{\E[f(X)]}{f(\epsilon)}.
}
\end{block}

\begin{Proof}<2->
Any monotonically non-decreasing function $f: \R \to \R_+$ is Borel measurable. 
Hence,  $f(X)$ is a random variable for any random variable $X$. 
Therefore, 
\EQ{
f(X) = f(X)\SetIn{X \ge \epsilon} + f(X)\SetIn{X < \epsilon} \ge f(\epsilon)\SetIn{X \ge \epsilon}.
}
%The result follows from the monotonicity of expectations. 
\end{Proof}
\end{frame}

\begin{frame}{Inequalities}
\begin{cor}[Markov]<1->
Let $X$ be a non-negative random variable, then $P\set{X \ge \epsilon} \le \frac{\E X}{\epsilon}$ for all $\epsilon  > 0$.
\end{cor}
\begin{cor}[Chebychev]<2->
Let $X$ be a random variable with finite mean $m_1$ and variance $\sigma^2$, 
then 
\EQ{
P\set{\abs{X -m_1} \ge \epsilon} \le \frac{\sigma_2}{\epsilon^2},\text{ for all }\epsilon  > 0.
} 
\end{cor}
\begin{Proof}<3->
Apply the Markov's inequality for random variable $Y = \abs{X-m_1} \ge 0$ and increasing function $f(x) = x^2$ for $x \ge 0$. 
\end{Proof}
\end{frame}

\begin{frame}{Inequalities}
\begin{cor}[Chernoff]<1->
Let $X$ be a random variable with finite $\E[e^{\theta X}]$ for some $\theta>0$, 
then 
\EQ{
P\set{X \ge \epsilon} \le e^{-\theta \epsilon}\E[e^{\theta X}],\text{ for all }\epsilon  > 0.
} 
\end{cor}
\begin{Proof}<2->
Apply the Markov's inequality for random variable $X$ and increasing function $f(x) = e^{\theta x} > 0$ for $\theta > 0$. 
\end{Proof}
\end{frame}

\begin{frame}{Inequalities}
\begin{Definition}[Convex function]<1->
A real-valued function $f: \R \to \R$ is convex if for all $x, y\in \R$ and $\theta \in [0,1]$, we have 
\EQ{
f(\theta x + (1-\theta)y) \le \theta f(x) + (1-\theta)f(y).
}
\end{Definition}
\begin{block}{Theorem : Jensen's inequality}<2->
For any convex function $f: \R \to \R$ and random variable $X$, we have 
\EQ{
f(\E X) \le \E f(X).
}
\end{block}
\end{frame}

\begin{frame}{Inequalities}
\begin{Proof}[Proof: Jensen's inequality]
\begin{itemize}
\item <1-> For simple random variables $X: \Omega \to \sX$, 
we show this by induction on cardinality of alphabet $\sX$.  
\item <2-> The inequality is trivially true for $\abs{\sX}=1$. 
We assume that the inductive hypothesis is true for $\abs{\sX} = n$. 
\item <3-> Let $X \in \sX$, where $\abs{\sX} = n+1$. 
Denote $\sX = \set{x_1, \dots, x_{n+1}}$ with $p_i \triangleq P\set{X = x_i}$ for all $i \in [n+1]$. 
\item <4-> $(\frac{p_j}{1-p_1}: j \ge 2)$ is a probability mass function for some random variable $Y \in \sY = \set{x_2, \dots, x_{n+1}}$ with cardinality $n$. 
\end{itemize}
\end{Proof}
\end{frame}

\begin{frame}{Inequalities}
\begin{Proof}[Proof: Jensen's inequality (Continued)]
\begin{itemize}
\item <1-> By inductive hypothesis, we have 
\EQ{
f\left(\sum_{i=2}^{n+1}\frac{p_i}{1-p_1}x_i\right) 
= f (\E Y) 
\le \E f(Y) 
= \sum_{i=2}^{n+1}\frac{p_i}{1-p_1}f(x_i).
}
\item <2-> Applying the convexity of $f$ to $\theta = p_1, x = x_1, y = \sum_{i=2}^{n+1}\frac{p_i}{1-p_1}x_i$, we get 
\EQ{
f(\E X) 
%= f(\sum_{i=1}^{n+1}p_ix_i)
= f\Big(p_1x_1 + (1-p_1)\sum_{i=2}^{n+1}\frac{p_i}{1-p_1}x_i\Big)
\le p_1f(x_1) + (1-p_1)f\Big(\sum_{i=2}^{n+1}\frac{p_i}{1-p_1}x_i\Big).
}
\item <3-> From the inductive step, it follows that the RHS is upper bounded by $\E f(X)$, and the result follows.
\end{itemize}
\end{Proof}
\end{frame}

\begin{frame}{$L^p$ spaces}
\begin{block}{Theorem}<1-> For any real numbers $1 \le p \le q < \infty$, we have $L^q \subseteq L^p$.
\end{block}

\begin{Proof}<2->
\begin{itemize}
\item <2-> Let $1 \le p \le q < \infty$ and consider a convex function $g: \R_+ \to \R_+$ defined by $g(x) \triangleq x^{{q}/{p}}$ for all $x \in \R_+$. 
\item <3-> It follows that $g(\abs{X}^p) = \abs{X}^q$ and hence from the Jensen's inequality, we get
\EQ{
\norm{X}_p^{q} = g(\E\abs{X}^p) \le \E g(\abs{X}^p) = \norm{X}_q^q.
} 
\end{itemize}
\end{Proof}
\end{frame}
\end{document}
